\SourceLeaf{157}
\Needspace{6\baselineskip}\pdfbookmark[3]{II. Valuation of obligations towards present members}{witwen-II-2-II}
\par\medskip{\raggedright\noindent II.\quad Valuation of the present monetary value of the widows' fund's obligations towards its present members.\par}\smallskip
As with the widows, I have also given all the members of the widows' fund consecutive numbers in the order in which they joined. Following throughout the path outlined in the place cited, I must first consider the 42 married members and, to begin with (which forms, as it were, the core of the investigation), determine the present monetary value both of the contributions they are obliged to pay and of the pensions their present wives will receive if they outlive their husbands, according to the principles of the calculus of probabilities and the mortality tables. Here I bring together on one page, in a table, the results of a long labour.

\begin{center}
\scriptsize\setlength{\tabcolsep}{1pt}
\begin{tabular}{r|c|r|r|r|r||r|c|r|r|r|r}
&\shortstack{Present\\married\\members}&\multicolumn{2}{c|}{\shortstack{Interest rate\\\(3\tfrac{1}{2}\) per cent}}&\multicolumn{2}{c||}{\shortstack{Interest rate\\4 per cent}}&&\shortstack{Present\\married\\members}&\multicolumn{2}{c|}{\shortstack{Interest rate\\\(3\tfrac{1}{2}\) per cent}}&\multicolumn{2}{c}{\shortstack{Interest rate\\4 per cent}}\\\cline{3-6}\cline{9-12}
No.&&\multicolumn{2}{c|}{\shortstack{Present monetary\\value of}}&\multicolumn{2}{c||}{\shortstack{Present monetary\\value of}}&No.&&\multicolumn{2}{c|}{\shortstack{Present monetary\\value of}}&\multicolumn{2}{c}{\shortstack{Present monetary\\value of}}\\
&&Pensions&\shortstack{Contri-\\butions}&Pensions&\shortstack{Contri-\\butions}&&&Pensions&\shortstack{Contri-\\butions}&Pensions&\shortstack{Contri-\\butions}\\\hline
&&Thl.&Thl.&Thl.&Thl.&&&Thl.&Thl.&Thl.&Thl.\\
109&L.&1752&59&1634&58&174&H.&1175&120&1053&114\\
135&O.&1509&85&1386&83&175&R.&1154&142&1007&134\\
138&C.&1846&70&1706&68&177&W.&1117&131&991&124\\
139&U.&937&109&854&104&180&T.&1064&132&944&125\\
140&H.&1451&106&1310&101&183&V.&943&148&824&139\\
141&L.&1091&93&1004&90&184&B.&1023&146&893&137\\
147&R.&1441&105&1301&101&187&W.&917&134&817&127\\
149&G.&1717&101&1545&97&188&S.&867&134&774&127\\
151&O.&1051&130&933&123&189&H.&1138&129&1011&122\\
153&B.&865&109&790&104&190&H.&936&142&825&134\\
155&K.&925&114&840&109&191&D.&892&123&804&117\\
156&S.&853&114&776&109&192&R.&943&139&848&131\\
160&H.&1376&123&1221&117&193&G.&1047&148&907&140\\
161&B.&876&120&795&114&195&D.&981&140&866&132\\
164&W.&1359&122&1209&117&197&R.&916&154&791&145\\
167&S.&995&140&878&132&198&W.&1072&144&936&136\\
168&K.&1345&135&1173&128&199&F.&1127&127&1005&120\\
169&Z.&1118&134&988&127&200&L.&861&153&749&144\\
170&R.&1541&99&1397&96&201&W.&1073&145&935&137\\
171&C.&1193&105&1084&101&203&B.&932&127&835&121\\
173&F.&1120&127&998&121&204&E.&787&146&729&138\\\cline{6-12}
\multicolumn{6}{r}{Total}&\multicolumn{2}{c|}{Thaler}&47324&5203&42366&4945
\end{tabular}
\end{center}
\noindent\hspace*{\parindent}I increase the total value of the widows' pensions by a sixth of itself on account of the orphans' pensions, not so much because \textsc{Kritter} assumed this ratio in his opinion as because it follows very nearly from my own discussion of the experience of our widows' fund [p. 154]. The total value of the contributions is then deducted, giving the net value of the fund's obligation towards the 42 married members. Finally, an addition is made simply pro rata for the 9 members who are now unmarried. These calculations are then as follows:
\begin{center}
\footnotesize\setlength{\tabcolsep}{3pt}
\begin{tabular}{l r|r}
&\shortstack{Interest rate\\\(3\tfrac{1}{2}\) per cent}&\shortstack{Interest rate\\4 per cent}\\
\shortstack[l]{Widows' pensions of the 42\\married members}&47324 Thl.&42366 Thl.\\
Orphans' pensions . . . . . . . . . . . . . . . . . . . . . .&7887 ,, &7061 ,, \\[5pt]
Pensions . . . . . . . . . . . . . . . . . . . . . . . . . . .&55211 ,, &49427 ,, \\
Contributions . . . . . . . . . . . . . . . . . . . . . . . . . . . .&5203 ,, &4945
\end{tabular}
\end{center}
\SourceLeaf{158}
\begin{center}
\footnotesize\setlength{\tabcolsep}{3pt}
\begin{tabular}{l r|r}
&\shortstack{Interest rate\\\(3\tfrac{1}{2}\) per cent}&\shortstack{Interest rate\\4 per cent}\\
\shortstack[l]{Net obligation of the fund towards\\42 married members}&50008 Thl.&44482 Thl.\\
\shortstack[l]{On that basis, proportionately, towards\\9 unmarried members . . . .}&10716 ,, &9532 ,, \\[5pt]
\shortstack[l]{Total value under the second heading\\ . . . . . . . . . . . . . . . . . . . . . . . .}&60724 ,, &54014 ,,
\end{tabular}
\end{center}
\noindent\hspace*{\parindent}I have already cautioned in my opinion, mentioned several times, that such a valuation cannot be accorded quite the same reliability as the calculation of the obligations under the first heading. For several circumstances arise in our widows' fund whose effect cannot be calculated in advance at all, and cannot even be estimated in the absence, for earlier times, of records such as those mentioned on p. [145, bottom]. I identify the remarriage of members after their wives' deaths as the most influential of these circumstances. For the wives' dying before their husbands is already among the chances taken into account in the calculation, and so every remarriage is a newly added charge on the fund's account that lies outside the calculation. On the other side, any other departure of a married member, for example through an appointment elsewhere, benefits the fund's account as an unestimated relief. The fact that widowed members who do not remarry nevertheless usually continue paying their contributions for at least a time after their wives' deaths belongs to the same category, although it is of little significance. Finally, perhaps the calculation for the presently unmarried members based on the average result for the married ones is somewhat too high. The body of unmarried members at any given time will usually be so composed that one or other of them is likely never to marry at all. This is, however, at least partly offset by the fact that the later the marriage takes place, the more probably one must presume that the wife will survive and have a long widowhood.

Although it is impossible to estimate in figures beforehand the effect of these various circumstances, one may nevertheless regard it as probable that they largely compensate for one another.

